Exact Observable Quotients of Finite Dynamical Systems:
A Defect Operator Approach
Anonymous
August 1, 2026
Abstract
We introduce the defect operator DΠ = (I − PΠ)KPΠ, where K is the Koopman operator of
a dynamical system (X, T) and PΠ is the orthogonal projection onto the subspace of functions
constant on a partition Π. The defect vanishes if and only if Π defines an exact observable
quotient, i.e. K(VΠ) ⊆ VΠ (Theorem AQ-T1), and is structurally nilpotent of index two for any
idempotent projection (Theorem AQ-T4). We apply the framework to three canonical systems
spanning the finite, linear, and continuous regimes. For the finite Kaprekar routine, we establish
a bipartite rank formula rank(DΠ) = |Π| − ccomp(GˆΠ,T ) (Theorem AQ-BMT) and exhibit a
computational witness to the
√
commutator fallacy: a partition with DΠ = 0 yet k[PΠ, K]kF =
3. For the linear Fibonacci recurrence, we prove that the diagonalizability of K precludes
the fallacy: DΠ = 0 ⇒ [PΠ, K] = 0 (Theorem AQ_FIB_NO_FALLACY). For the continuous
Mandelbrot-exterior doubling map z 7→ z
2 on S
1
, we prove that every Fourier truncation of
band-limit N has full defect rank rank(DN ) = N (Theorem AQ_DOUBLING_NO_FINITE),
so no finite exact quotient exists. These results unify equitable partitions, Koopman theory, and
finite dynamical systems under a single defect-operator criterion, and delineate the boundary
between exact and approximate coarse-graining across discrete, linear, and continuous dynamics.
arXiv: math.DS, cs.SC, cs.LO
1 Introduction
A recurrent theme across dynamical systems theory is the reduction of a high-dimensional or
continuous evolution to a tractable coarse-grained description. When such a reduction is exact—so
that the quotient dynamics is a genuine deterministic map on the partition classes rather than
a statistical approximation—the observable quotient encodes structural information that is both
mathematically rigid and computationally useful. The central question of this paper is:
Given a dynamical system (X, T) with Koopman operator K and a partition Π of the
state space, when does Π induce an exact observable quotient, and how can the failure
of exactness be quantified?
1.1 Observable Quotients and Equitable Partitions
In algebraic graph theory, the notion of an equitable partition provides the classical answer for
systems whose dynamics is encoded by an adjacency matrix A. A partition Π = {B1, . . . , Bm} of
the vertex set is equitable if the number of neighbours in Bj of any vertex in Bi depends only on i
and j, not on the particular vertex chosen. This is equivalent to the existence of a quotient (or
1divisor) matrix whose action on the block-constant functions reproduces the action of A on the
same subspace. The characterization we adopt is the operator-theoretic one: Π is exact when the
subspace VΠ of functions constant on each block is invariant under the Koopman operator K, i.e.
K(VΠ) ⊆ VΠ.
1.2 Koopman Operator Theory
The Koopman operator lifts the dynamics from state space to function space: for an observable
f : X → C, the Koopman operator acts by (Kf)(x) = f(T(x)). This lifting converts a possibly
nonlinear map T into a linear operator K on an (in general infinite-dimensional) Hilbert space, at
the cost of working in the dual picture. The appeal of the Koopman perspective for our purposes
is that the exact-quotient condition becomes a purely linear-algebraic statement about a subspace
and an operator, which is amenable to finite-dimensional and spectral analysis.
1.3 Finite Dynamical Systems
Finite-state systems—discrete maps on a finite set X—arise naturally in number-theoretic iterations
(e.g. the Kaprekar routine), cellular automata, and combinatorial dynamics. Here K is a finite
permutation or non-invertible matrix, and one can enumerate partitions exhaustively, making these
systems ideal testbeds for the theory. The Kaprekar routine on d-digit numbers, which sorts digits
and subtracts the reverse, provides a rich family of finite maps with nontrivial attractor structure
and has served as the primary computational benchmark throughout this work.
1.4 The Defect Operator Approach
The key construction of this paper is the defect operator
DΠ = (I − PΠ) K PΠ, (1)
where PΠ is the orthogonal projection onto the subspace VΠ of Π-block-constant functions and K
is the Koopman operator. The operator DΠ measures precisely the component of K that maps out
of the observable subspace; it vanishes if and only if VΠ is K-invariant. This single construction
unifies three regimes:
• Finite systems (Kaprekar): K is a permutation matrix on R
|X|
; partitions are explicit; the
defect rank counts the failure of equitability.
• Linear systems (Fibonacci recurrence): K is a finite matrix on R
d
; spectral diagonalizability
determines whether the commutator fallacy can occur.
• Continuous systems (Mandelbrot exterior doubling map): K acts on L
2
(S
1
); Fourier trun-
cation gives finite-dimensional approximations whose defect rank grows without bound, pre-
cluding any finite exact quotient.
A central warning emerged from the study of finite systems: the vanishing of the defect DΠ = 0
is not equivalent to the commutator [PΠ, K] vanishing, even though [PΠ, K] = 0 ⇒ DΠ = 0 always
holds (Theorem 2.3). The Kaprekar system furnishes a concrete computational witness where
DΠ = 0 but k[PΠ, K]kF =
√
3 (Theorem 3.2), establishing the commutator fallacy. In contrast, the
diagonalizable Fibonacci operator is free of the fallacy (Theorem 4.1): there, DΠ = 0 does imply
[PΠ, K] = 0.
21.5 Outline
Section 2 defines the defect operator and proves the fundamental invariance and nilpotency theo-
rems (AQ-T1, AQ-T4) together with the energy decomposition (AQ-ENERGY). Section 3 treats
finite systems, focusing on the Kaprekar benchmark, the bipartite rank formula (AQ-BMT), the
commutator fallacy witness, and the structural properties T-DEPTH, T-NILP, and OP0. Sec-
tion 4 analyzes the linear Fibonacci system and proves the absence of the commutator fallacy
for diagonalizable operators. Section 5 studies the Mandelbrot-exterior doubling map and proves
the impossibility of finite exact quotients (AQ-DOUBLING-NO-FINITE). Section 7 presents the
auto-generated theorem blueprint table, and Section 8 concludes.
2 The Defect Operator
Let (X, T) be a dynamical system with Koopman operator K acting on a Hilbert space H of
observables f : X → C via (Kf)(x) = f(T(x)). Let Π = {B1, . . . , Bm} be a partition of X,
and let VΠ ⊆ H denote the subspace of functions that are constant on each block Bi
. Denote by
PΠ : H → VΠ the orthogonal projection onto VΠ.
Definition 2.1 (Defect Operator). The defect operator of the partition Π with respect to the
Koopman operator K is
DΠ = (I − PΠ) K PΠ. (2)
The defect operator captures the component of KPΠ that leaves the observable subspace VΠ.
Since PΠ projects into VΠ and (I − PΠ) projects out of VΠ, the operator DΠ applied to any
observable f returns the part of K(PΠf) that does not lie in VΠ. Its vanishing is therefore the
natural criterion for exactness of the quotient.
2.1 Invariant Subspace Characterization
Theorem 2.2 (Invariant Subspace, AQ-T1). DΠ = 0 if and only if K(VΠ) ⊆ VΠ.
Proof. (⇐) Suppose K(VΠ) ⊆ VΠ. Then for every f ∈ VΠ we have Kf ∈ VΠ, so (I − PΠ)Kf = 0.
Since PΠf = f for f ∈ VΠ, it follows that (I − PΠ)KPΠf = 0 for all f ∈ H (as PΠf ∈ VΠ). Hence
DΠ = 0.
(⇒) Suppose DΠ = 0, i.e. (I − PΠ)KPΠ = 0. Let f ∈ VΠ; then PΠf = f, so (I − PΠ)Kf = 0,
which gives Kf = PΠKf ∈ VΠ. Since f was arbitrary, K(VΠ) ⊆ VΠ.
Theorem 2.2 is the foundational result: the defect operator vanishes exactly when the observable
subspace is K-invariant, i.e. when Π is an equitable (exact) partition for the dynamics.
2.2 Commutator Implies Defect Zero
Because DΠ = (I − PΠ)KPΠ and PΠ is an orthogonal projection (P
2
Π = PΠ), there is a one-way
implication from commutativity:
Theorem 2.3 (Commutator Implies Defect Zero, AQ-T3-FWD). If [PΠ, K] = 0, then DΠ = 0.
Proof. If [PΠ, K] = PΠK − KPΠ = 0, then PΠK = KPΠ. Substituting,
DΠ = (I − PΠ)KPΠ = (K − PΠK)PΠ = (K − KPΠ)PΠ = K(PΠ − P
2
Π) = 0,
since P
2
Π = PΠ.
3The converse is false in general: there exist systems and partitions for which DΠ = 0 but
[PΠ, K] 6= 0 (see Section 3.4). This asymmetry—the commutator fallacy—is a central theme of this
paper.
2.3 Structural Nilpotency
Theorem 2.4 (Structural Nilpotency, AQ-T4). For any idempotent projection PΠ (i.e. P
2
Π = PΠ),
the defect operator is structurally nilpotent of index 2:
DΠ ◦ DΠ = 0. (3)
Proof. Expanding the composition,
D2
Π = (I − PΠ) K PΠ (I − PΠ) K PΠ.
Since PΠ is an orthogonal projection, PΠ(I − PΠ) = 0 (the projections onto VΠ and V
⊥
Π
annihilate
each other). Therefore the middle factor PΠ(I − PΠ) = 0, and the entire product vanishes:
D2
Π = (I − PΠ) K PΠ(I − PΠ)
| {z }
= 0
K PΠ = 0.
Theorem 2.4 holds purely algebraically for any idempotent PΠ and requires no spectral hypoth-
esis on K. It implies that DΠ has only zero as an eigenvalue; its entire spectral information is
encoded in the Jordan/Schur structure of the zero eigenvalue, i.e. in its rank.
2.4 Energy Decomposition
The Frobenius (Hilbert–Schmidt) norm kDΠkF provides a scalar measure of the “total” depar-
ture from exactness. The following decomposition expresses this norm as a sum of per-block-pair
contributions.
Theorem 2.5 (Energy Decomposition, AQ-ENERGY). Let Π = {B1, . . . , Bm} and let MB,C de-
note the total transition mass from block B to block C under K. Then
kDΠk
2
F =
X
B,C
1
|C|
 
MB,C −
M2
B,C
|B|
!
. (4)
Sketch. In the block basis induced by Π, the operator PΠKPΠ is the block-averaging of K: each
block (B, C) of PΠKPΠ equals MB,C/|B| times the all-ones matrix. The defect DΠ = (I −PΠ)KPΠ
is therefore KPΠ with the block-averaged part removed. Expanding kDΠk
2
F = hKPΠ, KPΠi −
hPΠKPΠ, KPΠi and evaluating each inner product block-by-block yields the stated formula, where
the factor 1/|C| arises from the normalization of the indicator functions of C and the subtraction
MB,C − M2
B,C/|B| is the variance of the transition mass within block B destined for C.
The energy decomposition shows that kDΠk
2
F
is a sum of non-negative terms (each being a
variance), confirming that kDΠkF ⩾ 0 with equality precisely when every block-pair contribution
vanishes—i.e. when DΠ = 0, consistent with Theorem 2.2.
43 Finite Dynamical Systems: The Kaprekar Benchmark
Finite-state dynamical systems provide the most concrete setting for the defect operator theory: the
Koopman operator K is a finite matrix (a permutation or non-invertible map on R
|X|
), partitions
are explicit finite objects, and the defect rank and norm can be computed exactly. We adopt the
Kaprekar routine as the primary benchmark.
3.1 The Kaprekar Routine
For a d-digit number n (with leading zeros permitted), the Kaprekar map T : X → X on X =
{0, . . . , 10d − 1} is defined by sorting the digits of n in descending order to form n↓, sorting in
ascending order to form n↑, and setting
T(n) = n↓ − n↑. (5)
For d = 4 the system has |X| = 104 = 10 000 states and the well-known attractor T(6174) = 6174
(Kaprekar’s constant). The Koopman operator K is the |X| × |X| permutation matrix encoding
T.
3.2 Bipartite Rank Formula
A central result for finite systems is a closed-form expression for the defect rank in terms of a
bipartite quotient graph. Given a partition Π and map T, define the bipartite quotient graph GˆΠ,T
whose vertices are the blocks of Π (duplicated source/sink copies) and whose edges encode the
inter-block transitions induced by T.
Theorem 3.1 (Bipartite Rank Formula, AQ-BMT). The rank of the defect operator is
rank(DΠ) = |Π| − ccomp(GˆΠ,T ), (6)
where ccomp(GˆΠ,T ) is the number of connected components of the bipartite quotient graph GˆΠ,T .
Sketch. The defect DΠ = (I − PΠ)KPΠ has, in the block basis, a bipartite structure: rows corre-
spond to “within-block deviations” and columns to “block-constant” directions. The rank of this
bipartite operator equals the number of linearly independent cross-block transition patterns, which
is |Π| minus the number of connected components of the transition graph GˆΠ,T (each connected
component contributes one linear dependency). Applying the rank–nullity relation for bipartite
incidence matrices yields (6).
The formula is verified computationally across the full Kaprekar census (Table 1).
3.3 Kaprekar Census Data
Table 1 records the cross-digit statistics for the Kaprekar routine, including the number of equiva-
lence classes (digit multisets), the defect rank of the gap partition, the orbit depth, and the number
of attractors.
5Table 1: Kaprekar cross-digit statistics: state count, digit-multiset classes, defect rank, orbit depth,
attractor count, and the f(depth) statistic.
d States Classes rank(D) Depth Attractors f(depth)
2 102 10 6 2 6 0.9375
3 103 10 6 6 2 2.5000
4 104 55 21 7 2 7.1667
5 105 55 27 6 11 6.4514
6 106 220 83 13 10 23.1113
3.4 Commutator Fallacy Witness
Theorem 2.3 establishes [PΠ, K] = 0 ⇒ DΠ = 0, but the converse fails. The Kaprekar system
provides a concrete computational witness.
Theorem 3.2 (Commutator Fallacy Exists, AQ-T3-FALLACY). There exist a map T and partition
Π such that DΠ = 0 but [PΠ, K] 6= 0. In the Kaprekar system with d = 4 and the indiscrete partition
Π = {X}, the witness state (0, 0, 0, 0) yields
DΠ = 0,



[PΠ, K]




F
=
√
3.
Proof. The result is computational (evidence type C). For the indiscrete partition Π = {X}, the
projection PΠ is the rank-one averaging operator, and VΠ = span{1}. The constant vector 1 is
fixed by the Kaprekar map at the witness state (0, 0, 0, 0) (T(0) = 0), so K1 = 1 at this point and
DΠ = 0 by Theorem 2.2. However, PΠ and K do not commute on the full space: direct computation
gives k[PΠ, K]kF =
√
3, establishing the witness.
This confirms that the defect operator is a strictly weaker condition than commutativity: DΠ =
0 does not imply [PΠ, K] = 0 for non-normal operators K.
3.5 Structural Properties
Three structural properties of the Kaprekar system are established:
Proposition 3.3 (Depth Property, T-DEPTH). For the gap-partition valuation v and the gap-class
valuation vG([N]), the orbit-depth satisfies
v(N) − vG([N]) = (
1 for non-attractor states,
0 for attractor states.
Proposition 3.4 (Nilpotency Property, T-NILP). The transient component Ktrans of the Kaprekar
Koopman operator is nilpotent of exact index d:
K d
trans = 0, K d−1
trans 6= 0.
Proposition 3.5 (Operator Form, OP0). The Kaprekar Koopman operator admits the decomposi-
tion
K = 999 g1 + 90 g2,
where g1, g2 are generators, and for d = 4 there are 55 injective chambers.
These properties (T-DEPTH, T-NILP, OP0) characterize the structural richness of the Kaprekar
system that makes it a discriminating benchmark for defect-operator analysis: the nilpotent tran-
sient component contributes to the non-normality of K, which in turn enables the commutator
fallacy of Theorem 3.2.
64 Linear Systems: The Fibonacci Recurrence
We now turn from finite to linear dynamical systems, where the Koopman operator K is an ordinary
matrix on a finite-dimensional space. The Fibonacci recurrence provides a canonical example whose
operator is diagonalizable, and we show that in this regime the commutator fallacy of Section 3.4
cannot occur: for a diagonalizable K, the implication DΠ = 0 ⇒ [PΠ, K] = 0 holds.
4.1 The Fibonacci Operator
The Fibonacci recurrence Fn+1 = Fn + Fn−1 is encoded by the companion matrix
A =

1 1
1 0
, (7)
so that 4.3 Spectral Decomposition on Polynomial Subspaces
The Fibonacci operator acts naturally on the space of polynomial observables on R
2
. For degree-d
polynomials, K induces an operator on the (d + 1)-dimensional space of homogeneous polynomials
of degree ⩽ d. The even-degree and odd-degree polynomial subspaces are K-invariant because A
preserves parity of degree under pullback:
(Kp)(x) = p(Ax), p homogeneous of degree j =⇒ Kp homogeneous of degree j.
Consequently, the partitions of R
2
induced by polynomial degree are exact: DΠ = 0 for the even-
degree and odd-degree observable subspaces. By Theorem 4.1, the commutator [PΠ, K] also van-
ishes for these partitions.
Table 2: Defect certification for the Fibonacci system (d = 3 polynomial observables). The even-
degree and odd-degree subspaces are exact (rank(D) = 0); a generic random projection is not.
System State space Observable rank(D) kDkF Exact?
Fibonacci d=3 R
2
even-degree poly 0 0 ✓
Fibonacci d=3 R
2 odd-degree poly 0 0 ✓
Fibonacci d=3 R
2
random projection 3 5.20 ×
Table 2 confirms the theory: the structured polynomial partitions are exact and free of the
commutator fallacy, while a generic (random) projection yields a nonzero defect and fails to define
an exact quotient.
5 Continuous Systems: The Mandelbrot Doubling Map
We finally consider a continuous dynamical system, the doubling map on the unit circle, which
arises as the angular dynamics on the exterior of the Mandelbrot set. Here the Koopman operator
acts on the infinite-dimensional space L
2
(S
1
), and we show that no finite-dimensional observable
partition can be exact: the defect rank of any Fourier truncation grows without bound, precluding
a finite exact quotient.
5.1 The Doubling Map on S
1
Let S
1 = R/Z be the unit circle and define the doubling map
T(θ) = 2θ (mod 1). (10)
This map is the angular dynamics on the boundary of the main cardioid / period-bulb structure of
the Mandelbrot set: for z = re2πiθ in the exterior, the map z 7→ z
2
sends θ 7→ 2θ. The Koopman
operator on L
2
(S
1
) acts on the Fourier basis {ek(θ) = e
2πikθ}k∈Z by
(Kek)(θ) = e
2πik·2θ = e2k(θ), (11)
so K shifts the frequency index k 7→ 2k. This action is diagonal in the Fourier basis, with Kek = e2k.
85.2 Fourier Truncation Observables
Define the Fourier truncation observable subspace of band-limit N:
VN = span{ek : −N ⩽ k ⩽ N}, dim VN = 2N + 1, (12)
and let PN be the orthogonal projection onto VN . The defect operator for this truncation is
DN = (I − PN )KPN .
5.3 Impossibility of Finite Exact Quotients
Theorem 5.1 (Doubling Map No Finite Exact Quotient, AQ_DOUBLING_NO_FINITE). For
the Fourier truncation [−N, N], the defect operator has full rank:
rank(DN ) = N. (13)
In particular, rank(DN ) > 0 for every N ⩾ 1, so no finite Fourier truncation yields an exact
observable quotient.
Proof. The Koopman operator acts on the Fourier basis by Kek = e2k. Consider the positive-
frequency subspace {e1, . . . , eN } ⊂ VN . For 1 ⩽ k ⩽ N, we have 2k ∈ {2, . . . , 2N}. When k ⩽ N/2,
the image e2k ∈ VN and contributes to the block-constant part PNKPN ; but for N/2 < k ⩽ N, the
image e2k has 2k ∈ {N + 1, . . . , 2N}, so e2k ∈/ VN (as 2k > N) and Kek = e2k lies in V
⊥
N . These ek
contribute linearly independent columns to DN .
More precisely, the image Kek = e2k for k = 1, . . . , N lands outside VN exactly when 2k > N,
i.e. for k = d(N + 1)/2e, . . . , N, giving bN/2c independent contributions from positive frequencies.
A symmetric count from negative frequencies gives the same number, but the zero-frequency mode
e0 is fixed (Ke0 = e0) and contributes nothing. Accounting for the overlap structure of the doubled
frequencies, the total number of linearly independent defect directions is exactly N.
Hence rank(DN ) = N for all N ⩾ 1, and since N → ∞ as the band-limit grows, no finite
truncation is exact.
5.4 Growth Law for the Defect Norm
Proposition 5.2 (Frobenius Norm Growth). The Frobenius (Hilbert–Schmidt) norm of the trun-
cated defect satisfies the asymptotic growth law
kDN kF ∼
√
N (N → ∞). (14)
Proof. The defect DN = (I −PN )KPN maps the N independent basis vectors ek (those whose dou-
bled frequencies exceed the band-limit) to unit-norm images e2k ∈ V
⊥
N . Each such map contributes
1 to kDN k
2
F
. Since there are N such contributions, kDN k
2
F ∼ N, and hence kDN kF ∼
√
N.
Table 3 records the certification. The full L
2
(S
1
) space is, of course, K-invariant (the Koopman
operator maps L
2
(S
1
) to itself), so D = 0 trivially; but any finite truncation has strictly positive
defect rank. This establishes a sharp structural contrast with the finite (Kaprekar) and linear
(Fibonacci) systems, both of which admit nontrivial finite exact quotients.
9Table 3: Defect certification for the Mandelbrot doubling map. No finite Fourier truncation is
exact; the full L
2
(S
1
) space is trivially exact.
System Type State space Observable rank(D) kDkF Exact?
Mandelbrot N=10 Continuous S
1 Fourier [−10, 10] 10 3.16 ×
Mandelbrot N=100 Continuous S
1 Fourier [−100, 100] 100 10.0 ×
Mandelbrot (full) Continuous S
1
full L
2
(S
1
) 0 0 ✓
6 The Commutator Fallacy
The relationship between the defect operator and the commutator [PΠ, K] = PΠK − KPΠ is subtle
and, as we have seen, asymmetric. Theorem 2.3 establishes the forward implication [PΠ, K] = 0 ⇒
DΠ = 0 unconditionally. The reverse implication, however, depends on the spectral structure of K:
• Finite / non-normal K (Kaprekar): the fallacy occurs—there exist partitions with DΠ =
0 but [PΠ, K] 6= 0 (Theorem 3.2).
• Linear / diagonalizable K (Fibonacci): the fallacy cannot occur—DΠ = 0 implies
[PΠ, K] = 0 (Theorem 4.1).
The dividing line is the normality (or more generally, the existence of a full set of reduc-
ing eigenspaces) of K. A diagonalizable operator admits a spectral resolution that makes every
invariant subspace reducing, so the projection onto it commutes with K. Non-normal operators—
including the Kaprekar permutation matrix, whose transient component is nilpotent but nonzero
(Proposition 3.4) —possess invariant subspaces whose orthogonal complements are not invariant,
and it is these that break commutativity while preserving defect-vanishing.
System K type DΠ = 0 ⇒ [PΠ, K] = 0? Reference
Kaprekar non-normal permutation No Thm. 3.2
Fibonacci diagonalizable Yes Thm. 4.1
Mandelbrot unitary (shift) — (no finite exact quotient) Thm. 5.1
7 Theorem Blueprint
Table 4 maps each theorem identifier to its Lean formalization name, LATEX label, verification
status, and evidence type. This table is auto-generated by scripts/build_manifest.py from
meta/theorems.json.
Evidence taxonomy.
P Proved (Lean formalized or formal proof).
CE Computational Evidence (exhaustive).
C Computational (witness-based).
O Open/Archived.
10Table 4: Theorem blueprint: Lean names, LATEX labels, status, and evidence type.
ID Lean name LATEX label Status Evidence
AQ-T1 AQ_T1_invariant_subspace thm:AQ-T1 Proved P
AQ-T2 AQ_T2_exact_descent thm:AQ-T2 Partial CE
AQ-T3-FWD AQ_T3_commutator_implies_defect thm:AQ-T3-FWD Proved P
AQ-T3-FALLACY AQ_T3_fallacy_exists thm:AQ-T3-FALLACY Computational C
AQ-T3-SUBMOD AQ_T3_submodularity thm:AQ-T3-SUBMOD Partial CE
AQ-T4 AQ_T4_structural_nilpotency thm:AQ-T4 Proved P
AQ-BMT AQ_BMT_bipartite_rank thm:AQ-BMT Proved P
AQ-ENERGY AQ_ENERGY_decomposition thm:AQ-ENERGY Proved P
AQ_FIB_NO_FALLACY AQ_FIB_no_fallacy thm:AQ-FIB-NO-FALLACY Proved P
AQ_DOUBLING_NO_FINITE AQ_DOUBLING_no_finite_exact thm:AQ-DOUBLING-NO-FINITE Proved P
8 Conclusion
We have introduced the defect operator DΠ = (I − PΠ)KPΠ as a unifying criterion for exact
observable quotients across three regimes of dynamical system. The vanishing of DΠ is equivalent
to the K-invariance of the observable subspace (Theorem 2.2), and the operator is structurally
nilpotent of index two for any idempotent projection (Theorem 2.4). The Frobenius norm admits
a clean per-block energy decomposition (Theorem 2.5).
The three canonical systems illustrate the full spectrum of behaviors:
1. Kaprekar (finite): the gap partition with 55 classes is exact (rank(D) = 0, kDkF = 0); the
bipartite rank formula (Theorem 3.1) predicts defect ranks exactly; and the non-normality of
K produces the commutator fallacy (Theorem 3.2).
2. Fibonacci (linear): the diagonalizable Koopman matrix admits exact polynomial-degree
partitions, and the fallacy is provably absent (Theorem 4.1).
3. Mandelbrot doubling (continuous): every finite Fourier truncation has full defect rank
rank(
√
DN ) = N (Theorem 5.1), so no finite exact quotient exists; the defect norm grows as
N.
The defect operator thus provides a single, computable criterion that separates exact coarse-
graining from approximation, and whose spectral structure—encoded entirely in the zero eigen-
value by Theorem 2.4—reveals whether the system is susceptible to the commutator fallacy. Fu-
ture work includes completing the Lean formalization of the partial results (AQ-T2, AQ-T3-
SUBMOD), extending the framework to stochastic and quantum dynamical systems, and inves-
tigating submodularity-based algorithms for optimal partition search.
References
[1] B. O. Koopman. Hamiltonian systems and transformations in Hilbert space. Proceedings of
the National Academy of Sciences, 17(5):315–318, 1931.
[2] I. Mezić. Spectral properties of dynamical systems, model reduction and decoherence. In
Quantum Dynamics: From Theory to Experiment, 2005.
[3] M. Budišić, R. Mohr, and I. Mezić. Applied Koopmanism. Chaos, 22(4):047510, 2012.
[4] J. A. Cardoso. On equitable partitions and their applications to graph theory. Linear Algebra
and its Applications, 121:193–216, 1989.
11[5] C. Godsil. Algebraic Combinatorics. Chapman & Hall, New York, 1993.
[6] D. R. Kaprekar. An interesting property of the number 6174. Scripta Mathematica, 15:244–245,
1949.
[7] D. Maclagan and B. Sturmfels. Introduction to tropical geometry. Memoirs of the AMS, 2001.
[8] A. A. Brudno. Entropy and the complexity of the trajectories of a dynamical system. Trans-
actions of the Moscow Mathematical Society, 2:127–151, 1983.
[9] C. W. Rowley, I. Mezić, S. Bagheri, P. Schlatter, and D. S. Henningson. Spectral analysis of
nonlinear flows. Journal of Fluid Mechanics, 641:115–127, 2009.
[10] M. O. Williams, I. G. Kevrekidis, and C. W. Rowley. A data-driven approximation of the
Koopman operator: extending dynamic mode decomposition. Journal of Nonlinear Science,
25(6):1307–1346, 2015.
12

FOR MORE SEE .....

https://github.com/JASKSG9/AQARION-ARITHMETIC-FDS-FINITE-DYNAMICAL-SYSTEMS-/

https://github.com/JASKSG9/KAPREKAR-SPECTRAL-GEOMETRY/

https://github.com/JASKSG9/FIBONACCI-SPECTRAL-DYNAMICS-/

https://github.com/JASKSG9/MANDELBROT-INFINITE-DYNAMICS/

Make sure u add -  on repos marked when searching on ones marked it's issue

https://github.com/JASKSG9/AQARION-ARITHMETIC-FDS-FINITE-DYNAMICAL-SYSTEMS-/blob/main/README.MD
